% =====================================================================
\section{Primitives, Indexed Rendering and Memory}\label{sec:prims}
Every primitive is generated as a list of unique vertices $(x,y,z,n_x,n_y,n_z)$, 24 bytes each, plus a triangle
index list. Triangles are wound counter-clockwise seen from outside, so with \code{glFrontFace(GL\_CCW)} and
back-face culling only outward faces are rasterised. The angular resolution is $N=36$ slices,
$\Delta\phi=2\pi/36=10^\circ$, as written on the lab sheet. The generators are in \code{src/primitives/} and are
the \emph{only} source of geometry; the one exception, the terrain disc, is the cylinder cap subdivided into
rings and lifted by a height field (documented as the single non-rigid surface).

\begin{figure}[h]
  \centering
  \pictwo{primitive_library}{\linewidth}
  \caption{The primitive library rendered by the program (v2 capture, still valid in v4: the generators did not
  change). The counts are computed from the generators at capture time.}
  \label{fig:library}
\end{figure}

\subsection{Cube, sphere, cylinder, cone and frustum}
\textbf{Cube.} Each face has an outward normal $\vect n$ and in-plane axes $\vect u,\vect v$ with
$\vect u\times\vect v=\vect n$; its corners are
$\vect p_c=\tfrac12\vect n+\tfrac12 s_u\vect u+\tfrac12 s_v\vect v$ with
$(s_u,s_v)\in\{(-,-),(+,-),(+,+),(-,+)\}$, counter-clockwise about $\vect n$, split into $(0,1,2)$ and
$(0,2,3)$. A corner is stored once per face because its normal differs on each face: 24 vertices, 36 indices,
12 triangles (T01--T03).

\textbf{Sphere.} With $r=0.5$, $\theta_k=k\pi/S$, $\phi_j=2\pi j/N$ ($S=18$, $N=36$):
$x=r\sin\theta_k\cos\phi_j$, $y=r\cos\theta_k$, $z=r\sin\theta_k\sin\phi_j$, $\vect n=\vect p/r$. Only the
northern hemisphere is computed (a fan around the pole plus quads between rings). The southern half is its
reflection $y\mapsto-y$, as sketched on the lab sheet. A reflection reverses orientation, so each mirrored triangle
$(a,b,c)$ is emitted as $(a',c',b')$. The shared equator gives $2+(S-1)N=614$ vertices and $2N(S-1)=1224$
triangles with half the trigonometry (T04--T06).

\textbf{Cylinder, cone, frustum.} The axis is the local $z$ axis, $z\in[-0.5,0.5]$. The side joins two rings,
\begin{equation}
  \vect p^{\,b}_j=(r_b\cos\phi_j,\,r_b\sin\phi_j,\,-0.5),\quad
  \vect p^{\,t}_j=(r_t\cos\phi_j,\,r_t\sin\phi_j,\,+0.5),\quad
  \vect n_j=\mathrm{normalize}(\cos\phi_j,\,\sin\phi_j,\,r_b-r_t),
\end{equation}
and each cap is a fan around its own centre. With $r_b=r_t=0.5$ this is the lab cylinder: $4N+2=146$ vertices,
$4N=144$ triangles. With $r_t=0$ the same routine yields the \textbf{cone} (109 vertices, 72 triangles; the collapsed
top ring is skipped, so there are no degenerate triangles; T36); with $0<r_t<r_b$ and no caps it yields the
searchlight \textbf{frustum} (72 vertices, T10). Cones are used for pine whorls and the UAV missile noses (259 cone instances in the export), so the project still uses
exactly one cylinder generator for cylinder, cone and frustum.

\subsection{Indexed rendering and memory}
Without indexing every triangle stores its three vertices (lab ``technique 1'': 18 floats for a square); with
indexing unique vertices go to the VBO and triangles reference them through the EBO (technique 2: 12 floats $+$ 6
indices). The VAO records \emph{location 0: offset 0, stride 24\,B} (position) and \emph{location 1: offset 12\,B}
(normal); one VAO+VBO+EBO per primitive is shared by every part. Memory per mesh is
$M_\text{flat}=3T\cdot 24$\,B and $M_\text{indexed}=24V+12T$\,B (Table~\ref{tab:memory}, generated from
\code{scene\_export.json}). The six meshes take \memIndexed\,B indexed against \memNonIndexed\,B non-indexed; the
terrain disc (49\,921 vertices, 36\,m radius) dominates both. The mesh instance counts show how few distinct meshes
serve the \nLeaves{} parts. Test T12 checks that every indexed mesh is smaller than its flat version.

\begin{table}[h]
\centering\small
\caption{GPU meshes (rows from \code{build/scene\_export.json}, byte sums computed by \code{gen\_tables.py}). The
terrain's function string in the export still reads radius 20, a stale label in the exporter; the vertex count
(49\,921) is that of the 36\,m mesh (T13).}
\label{tab:memory}
% generated by report/gen_tables.py from build/scene_export.json -- do not edit
\begin{tabular}{@{}lrrrrrrr@{}}
\toprule
Mesh & $V$ & Indices & Triangles & Non-indexed (B) & Indexed (B) & Saving & Instances\\\midrule
Cube & 24 & 36 & 12 & 864 & 720 & 16.7\,\% & 1,177\\
Sphere & 614 & 3,672 & 1,224 & 88,128 & 29,424 & 66.6\,\% & 343\\
Cylinder & 146 & 432 & 144 & 10,368 & 5,232 & 49.5\,\% & 1,201\\
Cylinder (frustum) & 72 & 216 & 72 & 5,184 & 2,592 & 50.0\,\% & 2\\
Terrain disc & 49,921 & 292,320 & 97,440 & 7,015,680 & 2,367,384 & 66.3\,\% & 1\\
Cone & 109 & 216 & 72 & 5,184 & 3,480 & 32.9\,\% & 259\\
\midrule
Sum & & & & 7,125,408 & 2,408,832 & 66.2\,\% &\\
\bottomrule
\end{tabular}

\end{table}

\paragraph{Instancing and culling.} The grass is the unit cube drawn with \code{glDrawElementsInstanced} (one VAO,
a per-instance buffer with \code{glVertexAttribDivisor}); the preset decides the count: LOW 14\,000, HIGH 35\,000,
ULTRA 140\,000, EPIC 450\,000 blades (\code{qualityProfile}). Rain streaks and foliage cards use the same mechanism,
and the v4 effect sprites reuse the cube's $+Z$ face (Section~\ref{sec:pipeline}). Every pass (main view, UAV camera,
shadow maps, pre-pass) culls parts whose bounding sphere $r=\tfrac12\sqrt{s_x^2+s_y^2+s_z^2}$ lies outside one of the six
frustum planes of the view-projection matrix (Gribb--Hartmann).

\paragraph{Tests.} Unit tests T01--T13 check counts, centring, normal length, winding, sphere reflection, cylinder
axis, index validity and the memory comparison. Section~\ref{sec:testing} lists all \nTests{} tests.
